\chapter{Formal Specification of Rego}
\label{ch:form}

This chapter provides formal specification of the Rego policy language used in OPA.
These definitions are derived from the official OPA reference \cite{opa:web} and the OPA source code \cite{opa:git}.
All examples of Rego policies are valid code, interpretable online at the official Rego Playground: \url{https://play.openpolicyagent.org}.

%%%%%%%%%%%%%%%%%%%%%%%%%%%%%%%%%%%%%%%%%%%%%%%%%%%%%%%%%%%%%%%
% 80%
% variables, atomic values, composite values
\section{Terms and Values}
\label{s:terms}

%%===

This section establishes the vocabulary of values, variables, and terms that \rego policies operate over.
We define a substitution $\sigma$ mapping terms to their values; the rest of the chapter is built on this foundation.

%%===

\subsection{Variables and Substitution}
\label{ss:vars}

%%---

A \emph{variable} is an identifier matching the regular expression $[\_a-zA-Z][\_a-zA-Z0-9]^*$.
The symbol $\_$ alone is reserved as a placeholder for an \emph{unnamed variable}: a variable that is intentionally left unbound.
We denote the set of all variable names with $\Vars$.

Variables in a \rego policy are either \emph{global} or \emph{local}.
Global variables are accessible throughout the whole policy and their values compose the response document.
Local variables are scoped to the rule in which they are declared.
We write $\Globals \subset \Vars$ for the set of global variable names and $\Locals \subset \Vars$ for the set of local variable names, with $\Globals \cap \Locals = \emptyset$.
The internal OPA compiler renames all local variables to unique names prior to evaluation, so name shadowing does not arise in the formalization.

We define the set of \rego \emph{terms} $\Terms$ to encompass literals, variables, references, and function calls.
The central semantic object of this chapter is the \emph{substitution}
$$\sigma : \Terms \to \Vals$$
mapping each term to its value.
The codomain $\Vals$ includes the special element $\Undef$, whose properties are defined in Section~\ref{ss:ref}.
We call a term $t$ \emph{ground} (with respect to $\sigma$) if every variable occurring in $t$ is bound under $\sigma$; we call $t$ \emph{defined} if $\sigma(t) \neq \Undef$.
Every defined term is ground, but a ground term may still be undefined — for example, when a reference accesses a missing key in a collection.

The restriction $\sigma_V = \sigma|_{\Globals}$ to global variable names is precisely the response document returned by OPA to the querying service.

%%===

\subsection{Values}
\label{ss:vals}

%%---

The set $\Vals$ of \rego values is defined by the grammar:
\begin{align*}
v \;::=\;\;
  & s \in \Str
    & \text{(string literal)} \\
  \mid\;\; & n \in \Int
    & \text{(number)} \\
  \mid\;\; & b \in \Bool
    & \text{(boolean)} \\
  \mid\;\; & \Null
    & \text{(null)} \\
  \mid\;\; & [v_0, \ldots, v_{n-1}]
    & \text{(array)} \\
  \mid\;\; & \{v_{k_1}{:}v_1, \ldots, v_{k_n}{:}v_n\}
    & \text{(object)} \\
  \mid\;\; & \{v_1, \ldots, v_n\}
    & \text{(set)} \\
  \mid\;\; & \Undef
    & \text{(undefined)}
\end{align*}

Values of types $\Str, \Int, \Bool$, and $\Null$ are called \emph{atomic}; arrays, objects, and sets are called \emph{composite}.
$\Undef$ is a distinguished element that does not belong to any of the above domains; its formal properties are given in Section~\ref{ss:ref}.
The empty array is written $[\,]$, the empty object as $\{\}$, and the empty set as $\texttt{set()}$.
All composite types in \rego are \emph{heterogeneous}: their elements and field values may be of different types.

A \emph{literal} is a term whose value is determined syntactically.
The substitution maps every literal to itself: $\sigma(\ell) = \ell$.
A composite literal is evaluated element-wise; if any constituent term evaluates to $\Undef$, the entire composite literal evaluates to $\Undef$.

A variable $x \in \Vars$ is \emph{bound} under $\sigma$ if $\sigma(x) \neq \Undef$ and \emph{unbound} otherwise.
Assignment of a value to a variable is written $x \mathrel{:=} t$ and yields $\sigma(x) = \sigma(t)$:
\[
\textbf{(Assign)}\quad
\inferrule{
  x \mathrel{:=} t
}{
  \sigma(x) = \sigma(t)
}
\]
\rego provides two syntactic forms for assignment: the assign operator \texttt{:=} and the unification operator \texttt{=}\footnote{
  The unification operator \texttt{=} only functions as an assignment for an undefined left-hand side and it behaves as an equality test otherwise.
}.
This work uses the assign form throughout, in line with the OPA reference~\cite{opa:web:ass}.

%%===

\subsection{References}
\label{ss:ref}

%%---

A \emph{reference} is a term of the form $t[k]$, denoting access into the composite value of $t \in \Terms$ at key or index $k \in \Terms$.
This definition allows references to be chained into terms of form $t[k_1][k_2]\cdots[k_n]$; for such a reference $r$, we define its \emph{base} $\RefBase(r) = t$ and its \emph{length} $|r| = n+1$ (counting in the base $t$).
As syntactic sugar, $t\texttt{.}k$ abbreviates $t[\texttt{"k"}]$ when $k$ is a string literal identifier, so \texttt{x.a.b.c} is identical to \texttt{x["a"]["b"]["c"]}.

The value of a reference is defined recursively:
\[
\sigma(t[k]) \;=\;
\begin{cases}
  \Undef & \text{if } \sigma(t) = \Undef \text{ or } \sigma(k) = \Undef, \\
  \sigma(t)\bigl(\sigma(k)\bigr) & \text{if } \sigma(k) \in \mathrm{dom}\bigl(\sigma(t)\bigr), \\
  \Undef & \text{otherwise,}
\end{cases}
\]
where $\sigma(t)(k)$ denotes the lookup of key $k$ in the value $\sigma(t)$ — array indexing for arrays and field access for objects.
$\mathrm{dom}(v)$ then represents the set of all keys defined in a composite value $v$.
A chained reference evaluates left-to-right; if any intermediate step yields $\Undef$, the entire reference is $\Undef$.

Two special roots exist in every evaluation context.
The identifier \texttt{input} refers to the input document supplied by the querying service.
The identifier \texttt{data} is the root of a hierarchical document that combines the service-provided data (for example, a role-to-permission mapping) with the namespaces of all loaded \rego packages.
For a global variable $x$ defined in package $\pkg$, the reference \texttt{data.$\pkg$.x} and the bare name $x$ denote the same value within the same package: $\sigma(\texttt{data.\textit{$\pkg$}.x}) = \sigma(x)$.

The element $\Undef$ has the following properties that distinguish it from all concrete values:
\begin{itemize}
  \item Equality between two undefined terms is itself undefined: $\Undef$ does not compare equal to $\Undef$.
  \item Any expression containing a subterm that evaluates to $\Undef$ evaluates to $\Undef$.
  \item $\Undef$ is distinct from $\Null$: $\Null$ is a concrete JSON value used to explicitly represent the absence of information (for example, passed as a function argument to signal that an optional parameter is not provided), whereas $\Undef$ means the term is entirely absent from the evaluation result.
  \item undefined global variables do not appear in the output response document.
\end{itemize}

%%===
%% TODO: maybe change set type to contain values, instead of types
\subsection{Type System}
\label{ss:types}

%%---

Each value in $\Vals$ has a corresponding \emph{type}.
Since \rego is dynamically typed, the type of every term is determined at runtime from the concrete value produced by $\sigma$.
The set of \rego types $\Types$ is defined by the grammar:
\begin{align*}
\ty \;::=\; & \Str \mid \Int \mid \Bool \mid \Null \mid \Undef \\
   \mid\;   & \Array{\ty_0,\ldots,\ty_{n-1}} \\
   \mid\;   & \RegoSet{\ty_1,\ldots,\ty_m} \\
   \mid\;   & \Obj{v_1{:}\ty_1,\ldots,v_n{:}\ty_n}
\end{align*}

The intended meaning of each construct:
\begin{itemize}
  \item Atomic types $\Str, \Int, \Bool, \Null$ classify the corresponding atomic values;
        $\Undef$ classifies the undefined element.
  \item $\Array{\ty_0,\ldots,\ty_{n-1}}$ is an \emph{array type} of length $n$,
        where $\ty_i = \typof(v_i)$ is the type of the element at index $i$.
  \item $\RegoSet{\ty_1,\ldots,\ty_m}$ is a \emph{set type} whose element-type list
        is itself unordered and duplicate-free: $\{\ty_1,\ldots,\ty_m\}$ is the set of
        distinct types that appear among the elements, so $\typof(\{v_1,\ldots,v_k\}) =
        \RegoSet{\typof(v_1),\ldots,\typof(v_k)}$ with duplicates removed.
  \item $\Obj{v_1{:}\ty_1,\ldots,v_n{:}\ty_n}$ is an \emph{object type}
        with keys $v_i \in \Vals$ typed $\ty_i$.
\end{itemize}

We write $\typof : \Terms \to \Types$ for the function mapping each term to its type,
with $\typof(t) = \Undef$ whenever $\sigma(t) = \Undef$.
The analysis of how $\typof$ is approximated from a policy without a concrete evaluation
context is the subject of Chapter~\ref{ch:typ}.

\subsection{Functions}
\label{ss:functions}

%%---

A \emph{function} of arity $n \in \Nat_0$ is a pair $f = (A, \phi)$ where:
\begin{itemize}
  \item $A = (A_1, \ldots, A_n)$ is a list of \emph{argument type sets}, also denoted as $\sig(f)$, with $A_i \subseteq \Types$ being the set of types acceptable at position $i$,
  \item $\phi : \Vals^n \to \Vals$ maps concrete argument values to a result value.
\end{itemize}
We denote the set of all functions with $\Funs$ and write $\Funs : \Terms \to \Funs$ for the mapping from function name terms to their definitions.

Function names are terms in $\Terms$.
Their value under $\sigma$ is $\Undef$, so they do not contribute to the response document.

The substitution $\sigma$ is extended to a function call term $f(t_1, \ldots, t_n)$ by:
\[
\sigma\bigl(f(t_1,\ldots,t_n)\bigr) =
\begin{cases}
  \Undef & \text{if } \typof(t_i) \notin A_i \text{ for } \sig(f) = (A_1,\ldots,A_n) \\
  \phi\bigl(\sigma(t_1),\ldots,\sigma(t_n)\bigr) & \text{otherwise,}
\end{cases}
\]
where $(A, \phi) = \Funs(f)$.
Intuitively, the first case assigns the call an undefined value when any argument has a type not admitted by the function signature $A$, and since no function admits $\Undef$ as an argument type, the first case also subsumes undefined argument propagation.

Infix operators such as $+$, $-$, $==$, or $<$ are syntactic sugar for function calls and are rewritten to prefix form by the OPA compiler.
For example, the expression \texttt{x < y + 1} would be rewritten into \texttt{lt(x,plus(y,1))}.
% Every function additionally supports an \emph{assign form}: the return value is captured by supplying a target variable as an extra final argument, so $x \mathrel{:=} t_1 + t_2$ is equivalent to $\texttt{plus}(t_1, t_2, x)$.
% The OPA compiler normalizes to the assign form internally; this work uses the standard form for readability.

\rego also permits user-defined functions, whose definitions are expressed as rules.
Their formal treatment is deferred to Section~\ref{ss:fun-rules}.

%%%%%%%%%%%%%%%%%%%%%%%%%%%%%%%%%%%%%%%%%%%%%%%%%%%%%%%%%%%%%%%
\section{Expressions}
\label{s:exprs}
%%=============================================================

Expressions compose a condition of a rule body; further specified in Section~\ref{ss:rules:body}.
An \emph{expression} is one of the following:
\begin{itemize}
  \item a \emph{condition}: a term $t$ — including function calls — evaluated for its truth value.
        A condition holds iff $\sigma(t) \notin \{\Undef, \bot\}$, denoted $\Holds(t)$.
        A condition may optionally be negated with \texttt{not}, in which case it holds iff $\sigma(t) \in \{\Undef, \bot\}$.
  \item a \emph{declaration}: an introduction of locally scoped variables (with respect to the current rule body), in one of the following forms:
    \begin{itemize}
      \item \emph{assignment} $x \mathrel{:=} t$: binds $x$ to $\sigma(t)$ and holds iff $\sigma(t) \neq \Undef$.
            This is the primary form of local variable introduction used in this work.
      \item \emph{existential declaration} \texttt{some} [$k$,] $x$ [\texttt{in} $t$]:
            without \texttt{in}, this expression introduces an unconstrained variable $x$ with a value determined by the remaining conditions in the body.
            With \texttt{in}, followed by a collection $t \in \Terms$, $x$ quantifies the elements of $\sigma(t)$ and (optionally) $k$ quantifies its keys (with respect to $x$).
            The enclosing body is evaluated once for each such binding.
      \item \emph{universal declaration} \texttt{every} [$k$,] $x$ \texttt{in} $t$ \texttt{\{} $e_1 \cdots e_n$ \texttt{\}}:
            must be bound to a collection $\sigma(t)$, with $x$ quantifying its elements and (optionally) $k$ quantifying its keys (with respect to $x$).
            This declaration holds iff $e_1 \cdots e_n$ holds for every element of $\sigma(t)$.
            Variables introduced by \texttt{every} are scoped to its body and do not appear in the enclosing rule.
    \end{itemize}
\end{itemize}

Additionally, an expression can introduce a local variable using an \emph{assign form} of a function call.
This form is only applicable when the function call is the entire right-hand side of an assignment $x \mathrel{:=} f(t_1, \ldots, t_n)$, rewritten as $f(t_1, \ldots, t_n, x)$ by the OPA compiler.

%%%%%%%%%%%%%%%%%%%%%%%%%%%%%%%%%%%%%%%%%%%%%%%%%%%%%%%%%%%%%%%
\section{Policies and Rules}
\label{s:rules}
%%===

A \rego \emph{policy} is a finite set of rules $\Rules$ defined over a fixed package $\pkg$.
Every rule in $\Rules$ specifies a condition under which a particular term receives a value.
This section formalizes the structure of rules, their semantics, and the conditions under which a policy is consistent.
%%===

\subsection{Policies and Evaluation Context}
\label{ss:policy}

%%---

An \emph{evaluation context} is a pair $(\mathtt{input}, \mathtt{data})$ of ground JSON values, corresponding to the two documents supplied to OPA at query time (Figure~\ref{fig:opa-query-graph}).
Evaluating a policy $\Rules$ in context $(\mathtt{input}, \mathtt{data})$ produces a substitution $\sigma : \Terms \to \Vals$ satisfying all rules in $\Rules$; the formal conditions on $\sigma$ are developed in the subsections that follow.

The \emph{response document} is the restriction $\sigma_V = \sigma|_{\Globals}$, mapping each global variable name to its value.
Global variables for which no rule produces a result are absent from $\sigma_V$ — their value is $\Undef$ — unless a default rule provides a fallback (Section~\ref{ss:default}).

For a term $t \in \Terms$, we write $P_t \subseteq \Rules$ for the set of all rules whose head refers to $t$.

%%===

\subsection{Rule Head and Body}
\label{ss:head-body}

%%---

A \emph{rule head} $h$ is a structure from which we extract the following components:
\begin{itemize}
  \item $\RuleRef(h) \in \Terms$: the \emph{reference} of the head, identifying the term the rule contributes to.
    For complete and function rules, this is a single variable term; for accumulating object rules, this is a chained reference (e.g., \texttt{y.a.b}) identifying the object being built.
  \item $\RuleKey(h) \in \Terms \cup \{\Undef\}$: the \emph{key} contributed by this rule to $\RuleRef(h)$.
  \item $\RuleArgs(h) = (a_1, \ldots, a_n)$: the \emph{argument list}, empty for non-function rules.
  \item $\RuleDef(h) \in \{\top,\bot\}$: a boolean indicator of a \emph{default definition}, also called a \emph{fallback}.
\end{itemize}

In addition, we define the \emph{name of head $h$} as $\RuleName(h) = \RefBase(\RuleRef(h))$, i.e., the first term (variable) in the head reference.
For example, consider a rule \texttt{x.a.b := 1} with head $h$: $\RuleRef(h) = \texttt{x.a.b}$ and $\RuleName(h) = \texttt{x}$.

The type of a rule is determined by its head accessors as follows:
\begin{itemize}
  \item $\RuleKey(h) = \Undef$, $\RuleArgs(h) = ()$, $|\RuleRef(h)| = 1$:
        a \emph{complete definition} — assigns a single value to the global variable $\RuleRef(h)$,
  \item $\RuleKey(h) \in \Terms$, $\RuleArgs(h) = ()$:
        a \emph{set accumulating definition} — contributes one element to the set $\RuleRef(h)$,
  \item $\RuleKey(h) = \Undef$, $\RuleArgs(h) = ()$, $|\RuleRef(h)| > 1$:
        an \emph{object accumulating definition} — assigns a value to a nested location within the object identified by the prefix of $\RuleRef(h)$,
  \item $\RuleArgs(h) \neq ()$: a \emph{function definition} — defines a mapping from $\RuleArgs(h)$ to a value,
  \item $\RuleDef(h) = \true$: a \emph{default rule} — provides a fallback value for $\RuleRef(h)$ when no other rule applies; in such case $\RuleKey(h) = \Undef$.
\end{itemize}
The cases for $\RuleKey(h)$ and $\RuleArgs(h)$ are mutually exclusive: a rule cannot simultaneously be a function definition and an accumulating definition.
Each type is treated in its own subsection below.

A \emph{rule body} $B$ is a pair $(B_v, B_E)$, where:
\begin{itemize}
  \item $B_v \in \Terms$ is the \emph{value} (also denoted as $\BodyVal(B)$) the body contributes to the head; when absent from the source, $B_v = \true$,
  \item $B_E = (e_1, \ldots, e_m)$ is an ordered list of expressions forming the body condition, usually denoted as $\BodyCond(B)$.
\end{itemize}

A body $B$ is \emph{satisfied}, written $\Holds(B)$, iff every expression in $\BodyCond(B)$ holds:
$$\Holds(B) \iff \forall_{e \in \BodyCond(B)} : \Holds(e)$$
When $\BodyCond(B)$ is empty, $B$ is satisfied unconditionally.

A \emph{rule} is a pair $(h, \mathbf{B})$ where $\mathbf{B} = (B_1, \ldots, B_k)$ is an ordered list of bodies connected by \texttt{else}.
Body $B_i$ is only evaluated if no predecessor $B_j$ with $j < i$ was satisfied, i.e.:
$$\Holds(B_i) \text{ is considered only if } \forall \ranges{j}{1}{i-1} : \neg\Holds(B_j)$$

For convenience purposes, we abuse the notation of rule head's components $\RuleRef,\RuleDef,\RuleArgs, and \RuleName$ to be applicable to rules, since each rule only contains a single head; formally said, for $r = (h,\mathbf{B})$ we define:
\begin{align*}
  \RuleRef(r) &= \RuleRef(h) \\
  \RuleDef(r) &= \RuleDef(h) \\
  \RuleArgs(r) &= \RuleArgs(h) \\
  \RuleName(r) &= \RuleName(h)
\end{align*}

Given head $h$, we define the set $\Rules_h \subseteq \Rules$ called \emph{rules of $h$} as the maximum set $\{(h',\mathbf{B})|\RuleName(h) = \RuleName(h')\}$.
Informally said, for a head $h$, rules of $h$ form a set of all rules with heads of the same name as $h$.

\begin{figure}
  \label{fig:ex:auth}
  \begin{lstlisting}[language=Rego,escapechar=|]
package authentication

has_role(user, role) if { |\label{ex:has_role}|
  role in data.roles[user]
}

allowed_actions contains action if { |\label{ex:allowed}|
  some role in data.roles[input.user]
  some action in data.role_actions[role]
}

audit.user   := input.user |\label{ex:audit:user}|
audit.action := lower(input.action) |\label{ex:audit:action}|

default allow := false |\label{ex:default}|

allow if { |\label{ex:allow}|
  audit.action in allowed_actions
} else if {
  has_role(audit.user, "admin")
}
\end{lstlisting}
  \caption{Example \rego authentication policy.}
\end{figure}

\begin{example}
  Consider the policy from Figure~\ref{fig:ex:auth}.
  The rules on lines~\ref{ex:audit:user} and~\ref{ex:audit:action} are object accumulating
  definitions: both have $\RuleKey(h) = \Undef$, $|\RuleRef(h)| = 2$, and $\RuleName(h) = \texttt{audit}$,
  so they collectively build the object $\sigma(\texttt{audit})$ from two unconditionally satisfied contributions.
  The rule on line~\ref{ex:allowed} is a set accumulating definition with
  $\RuleKey(h) = \texttt{action}$, contributing one element per satisfied binding of the quantified variable $\texttt{action}$.

  The rule on line~\ref{ex:has_role} is a function definition with
  $\RuleArgs(h) = (\texttt{user}, \texttt{role})$.
  The rule on line~\ref{ex:allow} is a complete definition containing bodies $B_1,B_2$ (separated by \texttt{else}), where $\BodyVal(B_1) = \BodyVal(B_2) = \true$.
  If both $B_1,B_2$ would not hold, the variable specified in the head -- \texttt{allow} -- would be assigned the fallback value \false as defined on line~\ref{ex:default}.
\end{example}

\subsection{Complete Definitions}
\label{ss:complete}

%%---

A rule $(h, \mathbf{B}) \in \Rules$ is a \emph{complete definition} when $\RuleKey(h) = \Undef$,
$\RuleArgs(h) = ()$, $|\RuleRef(h)| = 1$, and $\RuleDef(h) = \bot$.
It assigns a single value to the global variable $\RuleRef(h)$: the value $\sigma(B_{i,v})$
of the first body $B_i \in \mathbf{B}$ whose conditions are satisfied.
\[
\textbf{(CompRule)}\quad
\inferrule{
  (h, \mathbf{B}) \in \Rules \\
  \Holds(B_i) \qquad \forall_{j < i} : \neg\Holds(B_j)
}{
  \sigma\bigl(\RuleRef(h)\bigr) = \sigma(B_{i,v})
}
\]
where $h$ satisfied the mentioned constraints for a complete definition.
When no body in $\mathbf{B}$ is satisfied, this rule contributes no value to $\RuleRef(h)$;
another rule in $\Rules_h$ may still do so, or a default rule provides a fallback
(Section~\ref{ss:default}).
If multiple rules in $\Rules_h$ have their conditions satisfied with differing values,
the policy is inconsistent; this is treated in Section~\ref{ss:consistency}.

\begin{example}
  The rule on line~\ref{ex:allow} of Figure~\ref{fig:ex:auth} is a complete definition
  for \texttt{allow} with $\mathbf{B} = (B_1, B_2)$.
  $B_1$ checks whether the audited action is permitted; $B_2$ is the \texttt{else} clause
  delegating to \texttt{has\_role}.
  Since $|\mathbf{B}| > 1$, the \texttt{else} semantics apply: $B_2$ is only considered
  when $\neg\Holds(B_1)$.
\end{example}

%%===

\subsection{Accumulating Definitions}
\label{ss:accum}

%%---

Rather than assigning a single value, accumulating definitions each contribute one element
to a collection associated with $\RuleName(h)$.
The final value is assembled from the contributions of all rules in $\Rules_h$ whose
conditions are satisfied.
Accumulating definitions take two forms.

\paragraph{Set accumulation.}
A rule is a \emph{set accumulating definition} when $\RuleKey(h) \in \Terms$ and $\RuleArgs(h) = ()$.
The element $\sigma(\RuleKey(h))$ is contributed to the set $\sigma(\RuleRef(h))$:
\[
\textbf{(AccSet)}\quad
\inferrule{
  (h, \mathbf{B}) \in \Rules \\
  \Holds(B_i) \qquad \forall_{j < i} : \neg\Holds(B_j) \\
  \sigma\bigl(\RuleKey(h)\bigr) \neq \Undef
}{
  \sigma\bigl(\RuleKey(h)\bigr) \;\in\; \sigma\bigl(\RuleRef(h)\bigr)
}
\]
where $h$ satisfies the constraints for a set accumulating definition.
The value of $\sigma(\RuleRef(h))$ is the set of all contributions across $\Rules_h$.
Set accumulation is always conflict-free: distinct elements are collected and duplicate
contributions are idempotent.

\paragraph{Object accumulation.}
A rule is an \emph{object accumulating definition} when $\RuleKey(h) = \Undef$,
$\RuleArgs(h) = ()$, and $|\RuleRef(h)| > 1$.
The chained reference $\RuleRef(h)$ identifies a nested location: its prefix is the
object being assembled and its last term is the key at which the value is stored.
The value $\sigma(B_{i,v})$ is assigned to that location:
\[
\textbf{(AccObj)}\quad
\inferrule{
  (h, \mathbf{B}) \in \Rules \\
  \Holds(B_i) \qquad \forall_{j < i} : \neg\Holds(B_j)
}{
  \sigma\bigl(\RuleRef(h)\bigr) = \sigma(B_{i,v})
}
\]
where $h$ satisfies the constraints for an object accumulating definition.
The object identified by the prefix of $\RuleRef(h)$ is assembled from all contributions
across $\Rules_h$.
A conflict arises when two rules assign different values to the same nested location;
this is treated in Section~\ref{ss:consistency}.

\begin{example}
  In Figure~\ref{fig:ex:auth}, the rule on line~\ref{ex:allowed} is a set accumulating
  definition with $\RuleRef(h) = \texttt{allowed\_actions}$ and $\RuleKey(h) = \texttt{action}$.
  The body is evaluated once per binding of \texttt{role} and \texttt{action},
  contributing each resulting value of \texttt{action} as an element of
  $\sigma(\texttt{allowed\_actions})$.
  The rules on lines~\ref{ex:audit:user} and~\ref{ex:audit:action} are object accumulating
  definitions with $\RuleName(h) = \texttt{audit}$ and $|\RuleRef(h)| = 2$.
  Both bodies are unconditionally satisfied, contributing two key-value pairs that
  together form $\sigma(\texttt{audit})$.
\end{example}

\subsection{Function Definitions}
\label{ss:fun-rules}

%%---

A rule $(h, \mathbf{B}) \in \Rules$ is a \emph{function definition} when $\RuleArgs(h) = (a_1, \ldots, a_n)$
is non-empty.
The arguments $a_i \in \Terms$ act as patterns: when the function is called with concrete
values, each $a_i$ is unified with the corresponding argument, producing local bindings
under which the body is evaluated.
The semantics is otherwise identical to a complete definition — the first satisfied body
$B_i$ determines the return value, and the \texttt{else} chain applies.

User-defined functions enter $\Funs$ (as defined is Section~\ref{ss:functions}) through their rules: for a function name $f$
defined by rules, $\Funs(f)$ is constructed from the argument patterns $\RuleArgs(h)$
and the rule bodies, with the argument type set $A$ derived from the types of the patterns.
This connects function definitions back to the formal definition of functions in
Section~\ref{ss:functions}.

\begin{example}
  The rule on line~\ref{ex:has_role} of Figure~\ref{fig:ex:auth} is a function definition
  with $\RuleArgs(h) = (\texttt{user}, \texttt{role})$.
  When called as $\texttt{has\_role}(\texttt{audit.user},\, \texttt{"admin"})$,
  the arguments are unified with \texttt{user} and \texttt{role} respectively,
  and the body \texttt{role in data.roles[user]} is evaluated under those bindings.
\end{example}

%%===

\subsection{Default Rules}
\label{ss:default}

%%---

A rule $(h, \mathbf{B}) \in \Rules$ is a \emph{default rule} when $\RuleDef(h) = \top$.
It provides a fallback value for $\RuleRef(h)$ when no other rule in $\Rules_h$ contributes
a value.
The \rego language enforces default rules to contain a single body $(v,())$, where $v \in \Vals$, i.e., given value must be constant and there are no conditions for the body.
Furthermore, given any rule head $h$, at most one default rule is allowed in $\Rules_h$; formally said:
$$\forall r,r' \in \Rules_h: \RuleDef(r) \implies (\neg\RuleDef(r') \lor r = r')$$

We define $\Def : \Terms \to \Vals$ mapping each head reference to its default value,
with $\Def(t) = \Undef$ when no default rule exists for $t$.
The semantics is:
\[
\sigma\bigl(\RuleRef(h)\bigr) = \Def\bigl(\RuleRef(h)\bigr)
\quad \iff \quad
\forall\,(h', \mathbf{B}') \in \Rules_h \setminus \{(h,\mathbf{B})) :
\forall\, B'_i \in \mathbf{B}' : \neg\Holds(B'_i)
\]
Without a default rule, a head whose no other rule applies is $\Undef$ and absent from
the response document.
With a default rule, the head always carries a defined value.
Default rules are applicable only to complete definitions and are allowed at most once
per head.

\begin{example}
  The rule on line~\ref{ex:default} of Figure~\ref{fig:ex:auth} sets
  $\Def(\texttt{allow}) = \false$.
  If neither body of the \texttt{allow} rule on line~\ref{ex:allow} is satisfied,
  $\sigma(\texttt{allow}) = \false$ rather than $\Undef$, ensuring \texttt{allow}
  always appears in the response document.
\end{example}

%%===

\subsection{Consistency}
\label{ss:consistency}

%%---

A policy $\Rules$ is \emph{consistent} in a given evaluation context if no two rules
in $\Rules$ contribute differing defined values to the same location.
Let $v(h, \mathbf{B})$ denote the value contributed by the first satisfied body of
$(h, \mathbf{B})$, or $\Undef$ if none is satisfied.
Consistency then requires:
\begin{gather*}
  \forall\,(h_1,\mathbf{B}_1),(h_2,\mathbf{B}_2) \in \Rules,\;
  \RuleRef(h_1) = \RuleRef(h_2)\;: \\
  v(h_1,\mathbf{B}_1) = v(h_2,\mathbf{B}_2)
  \;\lor\;
  \Undef \in \bigl\{v(h_1,\mathbf{B}_1),\, v(h_2,\mathbf{B}_2)\bigr\}
\end{gather*}
Set accumulation is always conflict-free: multiple rules contributing the same element
are idempotent and distinct elements are simply collected.
When inconsistency is detected in the current evaluation context, OPA reports it as a
runtime error.

\begin{example}
  Consider two complete definitions for \texttt{allow}:
  \begin{lstlisting}[language=Rego]
allow := true  if { input.role == "editor" }
allow := false if { input.age < 18 }
\end{lstlisting}
  For an evaluation context where \texttt{input.role} is \texttt{"editor"} and \texttt{input.age} equals $18$, both bodies are satisfied with
  differing values $\true$ and $\false$, so $v(h_1, \mathbf{B}_1) \neq v(h_2, \mathbf{B}_2)$
  and neither is $\Undef$ -- a conflict.
  A solution to this conflict would be chaining the two bodies into a single rule, implying a precedence to the evaluation.
\end{example}
